\documentclass[../main]{subfiles}
\begin{document}
\ifSubfilesClassLoaded{\setcounter{page}{27}}{}
\section{Common (Social) Property}
\label{sec:05_common_property}

The obstacle to the movement towards communism is private property. The first work in which Marx gives serious and thorough consideration to communism is the “Economic and Philosophic Manuscripts of 1844”. The chapter on communism appears in the second manuscript, immediately following the chapter on private property. It opens with reflections on private property, on labour as the essence of private property, and on how private property and labour are understood by various representatives of socialism and communism as an ideological current. Communists (including those before Marx) oppose private property in the means of production because it makes people the slaves of other people and gives rise to all the horrors of capitalist society.

Private property and personal property are not the same thing. Items of direct use and consumption belonging to their owner are, in form, private property, just as plants, factories, shops, corporations and large tracts of land are. This circumstance obscures the essence of the matter. That is why, when anti-communists hear about the abolition of private property, the most primitive among them think it means the socialisation of absolutely everything and start fearing for their personal belongings. On this point, we can reassure them. Communists are interested in articles of personal consumption not in the sense of depriving anyone of them, but only in ensuring that everyone is provided with them to a degree sufficient for a full and many-sidedly rich life. Economically, these items are not private property but personal property.

The abolition of private property does not do away with personal property (private property is a relationship between people concerning things, not the things themselves), since consumption must necessarily be personal in character. On the contrary, the abolition of private property ensures that all members of society have access to the items they need for personal consumption, and, where it is more rational to do so, it shifts consumption into social forms. For example, organising consumer cooperatives makes it possible to use things more efficiently and reduces the need to produce and store them, and therefore also the amount of social time spent on this. Thus, organising a simple consumer cooperative for sharing items that are not used very often makes it possible to avoid cluttering up space and to use each item as efficiently as possible right up to its moral and physical depreciation. The item passes from hand to hand, and every member of the cooperative can derive a useful effect from it. Examples of consumer cooperatives have been described in the literature many times. This is easily done with a wide variety of items, which cease to belong to anyone personally while still being used by everyone to the extent needed to satisfy their needs. Here, communism stands opposed to personal property not at all from the side its opponents fear. Social forms of consumption are aimed not at taking something away from someone, but at giving everyone access to articles of consumption.

As for such social forms of everyday life as the shared use of washing machines, communal cooking and food delivery, and the shared use of cleaning equipment, the savings in resources and human time (both directly and indirectly—as savings in the time spent producing those resources) can be very large while maintaining the same level of consumption. At any given moment, these savings are proportional to the capabilities of the technology. But even now, for example, cooking porridge for one child at home takes the same amount of time as cooking porridge for 200 children in a kindergarten or a milk kitchen equipped with modern appliances.

The same applies to cooking for adults. Social consumption makes it possible to use more efficient technologies that cannot be applied in individual households, with a much higher level of automation in every process. At the same time, quality control can be carried out better, which ultimately benefits everyone involved. Movement in this direction exists under capitalism as well. Capitalism itself, in its own way, produces and standardises systems of public catering (McDonald's, for example, offers identical French fries in every part of the world), laundries, and so on. Under capitalism, however, they are still not used in a social way. Forms of collective consumption constantly run up against resistance in commodity production, and production, as we know, determines consumption.

All the burdens of everyday life, in which a great many people are engaged, can be replaced by advanced production based on modern technology; private transport, which gets in its own way and that of pedestrians and creates traffic jams, is replaced by a well-developed public transport system, and so on, in every sphere of human life where this would benefit everyone. For example, in today's large cities with developed public transport, those who use it, especially at rush hour, reach their destinations more efficiently. In Moscow, children who take the metro to school mostly arrive on time, whereas those who are driven in private cars tend to be late, because attempts to use limited public urban space in a private way prove less efficient than its social use. This also explains why city authorities, during snowfalls, urge people to give up their private cars and switch to public transport in order to minimise accidents on the roads.

Commodity production cultivates individual forms of consumption. For all their irrationality and inefficiency, they are a condition for constant demand for commodities on the market, a condition for the movement of capital. It is precisely in this sense that private property comes into contact with personal property and influences it.

Private property, or, to put it more precisely, private property in the means of production, is above all the ability to extract surplus product (surplus value, in the form of profit or in the form of rent), that is, the exploitation of other people.

Communists speak of the socialisation of private property, not personal property—of the means of social production, which today are used not for the benefit of society as a whole but serve the exploitation of man by man. Private property is not things, but social relations mediated by things. Its essence is exploitation.

Representatives of communism before Marx understood the abolition of private property as its extension to the whole of society, the transformation of all members of society into workers, and the equalisation of everyone with respect to both labour and consumption. Many representatives of pre-Marxian communism regarded private property merely as material wealth that could be owned by everyone in common. Proudhon, for example, conceived of private property only as capital, which was to be abolished. Closer to Marx were those theorists who saw the source of private property and of its pernicious influence on man in a particular character of labour, “labour which is levelled down, fragmented, and therefore unfree”, yet they reduced such labour to particular kinds of labour. “Fourier, who, like the Physiocrats, again regards agricultural labour as at least the best kind of labour, whereas for Saint-Simon, on the contrary, the essence of the matter lies in industrial labour as such, and he accordingly aspires to the exclusive rule of the industrialists and the improvement of the workers’ condition. Finally, communism is the positive expression of the transcendence of private property; at first it appears as universal private property.” “Positive” here means “preserving all the wealth of culture and the active capacities of man developed in a society of commodity production”.

Marx subjects this understanding of the negation of private property as its extension to the whole of society to merciless criticism. He shows that such communism does not overcome private property but is merely the theoretical culmination of private property as a social relation. He notes that the dominance of property, as an aggregate of things, over such communism, both as an ideological and theoretical position and as a social movement, is still very great. Social relations are still conceived only in terms of private property. Such communism “wants to destroy everything which is not capable of being possessed by all as private property. It wants to disregard talent, etc., in an arbitrary manner. For it the sole purpose of life and existence is direct, physical possession. The category of the worker is not done away with, but extended to all men. The relationship of private property persists as the relationship of the community to the world of things. Finally, this movement of opposing universal private property to private property finds expression in the brutish form of opposing to marriage (certainly a form of exclusive private property) the community of women, in which a woman becomes a piece of communal and common property.”

For Marx, the idea of the community of women is a relation to a human being as a thing, which also reveals man’s relation to the world of things created by man for man: “Just as woman passes from marriage to general prostitution, so the entire world of wealth (that is, of man’s objective substance) passes from the relationship of exclusive marriage with the private owner to a state of universal prostitution with the community.” Marx calls such communism crude. Wealth here means the entire aggregate of things created by people for people, in which are embodied both the active capacities of human society and the needs it has developed.

The theory of crude communism, which negates the human personality, merely absolutises in thought a real tendency of the society of private property. Where the wealth of human relations is reduced to relations of private property, the personality is reduced to an economic agent of private property. Indeed, the personality of these agents (the proletarian and the capitalist alike) is a matter of complete indifference to the movement of private property. The capitalist is interested in the proletarian not as a personality but solely as the seller of labour power of a certain quality, and the proletarian is interested in the capitalist solely as its buyer. The same thing happens in every act of buying and selling any other commodity. All the other personal qualities of economic agents have no bearing on the movement of private property. If people spend their entire lives without ever stepping outside these relations, their personality begins to be reduced to these determinations alone.

Crude communism presupposes a transcendence of private property that would not be a genuine assimilation of it. This is evident from its abstract negation of the achievements of culture, its treatment of them as superfluous merely because the proletarian, as a proletarian, has no need of them, and (in Marx’s words) its striving for the unnatural simplicity of the poor and crude man who has few needs, who cannot overcome private property because he has not yet even reached it.

The slogan “Take everything and divide it up”, typical of present-day exponents of pseudo-communism, stands lower still. They have not even reached socialisation in the form proposed by the exponents of crude communism, namely the socialisation of poverty.

Against the crude communist theories, Marx sets communism as the positive transcendence of private property. Such communism, as a historical process, preserves the wealth of all preceding historical development in the form of the active capacities, needs, and material and spiritual culture of humanity. Communism makes all of this the possession of every member of society. As a historical process, it is conscious rather than spontaneous in character, since socialisation presupposes a conscious attitude on the part of the members of society towards the production and reproduction of their social life. Communism presupposes the assimilation of the objective laws of social development, which humanity can put to its service in exactly the same way as it does the laws of physics, chemistry and biology in modern industry.

The whole preceding history of humanity creates the conditions for a communism that overcomes private property while preserving, in this very overcoming, all the wealth of preceding development. Over the course of all preceding history, a corresponding consciousness also takes shape: the theory of communism. Communism is understood as the process of the abolition of private property: “For the abolition of the idea of private property, the idea of communism is quite sufficient. For the abolition of private property in actual reality, however, actual communist action is required.”

One point deserves mention, as we consider it important. It concerns the use of the terms “socialism” and “communism”, whose familiar meanings did not become established at once. In the Manuscripts of 1844 quoted above, “socialism” still refers to a society in which private property is not in the process of being sublated but has already been sublated. Communism, by contrast, is understood as the process (act) of negating (sublating) private property. That is precisely why it is regarded here not as a form of society but merely as the dynamic principle of the immediate future.

This circumstance sheds light on why Marx and Engels raised precisely these questions, and set them out in precisely this way, in the “Manifesto of the Communist Party”, written just three years after the “Manuscripts”. The roots of the Manifesto’s central proposition become more clearly visible: “But modern bourgeois private property is the final and most complete expression of the system of producing and appropriating products that is based on class antagonisms, on the exploitation of the many by the few. In this sense, the theory of the Communists may be summed up in the single sentence: Abolition of private property.”

This sublation, as already noted, presupposes the preservation of all the cultural wealth, active capacities and needs developed on the basis of private property. It is carried out not for its own sake but for the existence of a communist society with its corresponding property and form of intercourse.

Whatever form communist society may take, man as a natural being must live from nature and appropriate nature: everything that man uses and consumes is made from natural material, by means of natural processes put to man’s service, or is itself those natural processes. Twenty years after writing the “Manifesto”, Marx wrote in the Introduction to the “Economic Manuscripts of 1857–1859”: “All production is appropriation of nature by an individual within and through a specific form of society. In this sense it is a tautology to say that property (appropriation) is a precondition of production. But it is ridiculous to leap from this to a definite form of property, e.g. private property (which, moreover, presupposes an antithetical form, non-property, as a condition).”

Communist society corresponds to a common, that is, social form of property in all the social means of production and, accordingly, to common social management of all social production as a single whole. The social character of production, developed on the basis of private property, must be matched by the social character of appropriation of the products of that production. Marx and Engels wrote about this dozens of times in various works. This is perhaps the only proposition that modern communists have remembered, even if they do not understand its significance.

To set this historical movement in motion in the most advanced countries, Marx and Engels proposed the following measures: “1. Abolition of property in land and application of all rents of land to public purposes.

2. A heavy progressive tax. 3. Abolition of the right of inheritance. 4. Confiscation of the property of all emigrants and rebels. 5. Centralisation of credit in the hands of the state, by means of a national bank with state capital and an exclusive monopoly. 6. Centralisation of all transport in the hands of the state. 7. Extension of factories and instruments of production owned by the state; the bringing into cultivation of land and the improvement of the soil in accordance with a common plan. 8. Equal liability of all to work; establishment of industrial armies, especially for agriculture. 9. Combination of agriculture with industry; promotion of the gradual elimination of the distinction between town and country. 10. Public and free upbringing of all children. Abolition of children’s factory labour in its present form.”

As we can see, given the stated aim of the abolition of private property, these measures do not look all that radical; rather, they are general, as befits a manifesto. Manifestos must be clear and comprehensible, and they must indicate the main direction of development. Moreover, the authors themselves maintained that the measures depend on concrete-historical conditions and may therefore differ. We will discuss these measures specifically in one of our upcoming issues, through the prism of the experience of 20th-century socialism.

For now, a few more words need to be said about what private property is, and what exactly communists set out to abolish while preserving all of its positive content, so that “in place of the old bourgeois society, with its classes and class antagonisms,” we shall have “an association, in which the free development of each is the condition for the free development of all.”

Private property is the product of unfree labour and, at the same time, a means of production of the unfree character of labour: labour in which neither the process nor the product of labour belongs directly to the collective of working people; labour that is, in itself, a matter of indifference to the workers and serves merely as a means of achieving external ends (money). The essence of private property lies in this character of labour. It reaches the limits of its development in capitalist society with its developed commodity production, thanks to the social division of labour being carried to its extreme, turning the products of labour into commodities (in commodity production).

For a long time, the social division of labour furthered the development of humanity. It raised the productivity of labour and created the material basis for transferring the machine-like functions that people are forced to perform because of underdeveloped tools to those very tools, to machines. But it does not merely unite people along the lines of “Vasya is a fitter and Petya is a programmer, each does his own job, and above them stands a specialist manager who directs them for the common good.” No participant in the process pursues such a goal. Under this state of affairs, all people end up connected to one another externally—through the products of their labour. They are forced to exchange their product for one they do not have, while remaining indifferent to one another in the labour process. It does not matter who makes this product or how; what matters is that I need this product of someone else’s labour.

The social division of labour has passed through several major stages that played a positive role in the development of humanity. They deserve to be discussed in a separate work devoted specifically to this question. Here it is only important to note that social development is a product of history. Marxists reject outright all prejudices about the biological or natural suitability of members of society for particular kinds of labour. These prejudices do not withstand even superficial scientific criticism. They are at odds not only with what is known about the history of this question, but also with the history unfolding before our very eyes. Every year a multitude of new professions appear and disappear. Parents often do not understand what their children’s work is actually about, and they certainly could not have passed on an aptitude for it through inheritance, along with their genes.

From the moment labour becomes partial and the products of people’s labour become commodities, that is, begin to be produced specifically for exchange, these products begin to dominate people. Things—the products of their own labour—and money, as the representative of all things, acquire a well-known power over living people. And from the moment labour power becomes a commodity, such domination is accordingly acquired by those who have enough of this money to buy means of production and labour power. These people own living human labour. Economically, they are merely representatives of capital. Their function is to attach other people’s labour to capital so that the value of capital increases. Ultimately, this leads to a situation in which almost all members of society, in one way or another, work for the sake of increasing capital—the self-expansion of value. “Only such products can become commodities with regard to each other, as result from different kinds of labour, each kind being carried on independently and for the account of private individuals.” Therefore the abolition of private property is the abolition of commodity production in general, and not only of capitalist commodity production.

Value strives for self-expansion. This is the law that governs a society in which everyone is chained to their own partial activity—in other words, a society of developed commodity production. The laws of modern states are the natural juridical expression of this objective law of production, which is based on unfree and fragmented labour.

Up to a certain point, capital's drive for self-expansion is the driving force of social development. But it calls into being social forces that become cramped within the confines of capital. Using them exclusively for the purpose of increasing capital leads to catastrophic consequences.

These forces must be managed—and indeed can only be managed for the good of society—by everyone together, for the development of all, according to a common rational plan. And for this it is necessary to master universal human capacities (to acquire universal sensuousness, thinking and morality), going beyond the confines of one's profession, which do not allow one to make sense even of a “neighbouring” profession, let alone of social processes as a whole. Therefore the abolition of private property is equivalent to the abolition of the unfree character of labour. It is accomplished as free communist activity develops and becomes the common possession of all people.

To private property, communism opposes common property as a form of people's joint practical relationship to the means of production and natural resources. Common (Social) Property ensures the most efficient use of productive capacities not only for society as a whole but also for each of its members. It presupposes active participation and an active will with respect to social resources on the part of every member of society. Accordingly, if private property is a product of the division of labour (its fragmentation into ever more particular kinds of labour)\footnote{It was precisely because the social division of labour in the USSR had not been overcome (the tendency towards overcoming it had not become predominant), and the proletariat was reproduced as a proletariat, that is, as a class of bourgeois society, that the complete restoration of the rule of private property took place.}, and is in fact this very division, then common property presupposes, as noted above, the socialisation of labour and the elimination of its division. Social property presupposes free human activity by the standards of truth, goodness and beauty. It is this that constitutes the real economic basis for a real, and not merely formal, socialisation, and for a real coincidence of the private and social interests of all members of society.

Common or social property should not be understood as the property of collective farms and cooperatives, or of other groups of people who produce commodities and trade them with the rest of the world. These do not go beyond the bounds of private property, since they do not go beyond the bounds of commodity production. The members of such enterprises are connected with society as commodity owners and commodity producers. The very nature of commodity production also determines the relations within these collectives, since it puts the production of value first and makes profit the goal of the collective's work.

Nor should the germ of common property be sought, for example, in the property of the family. In the 19th century, some socialists imagined that social property could arise through the “expansion” of family property to the scale of society as a whole. Society was pictured as a large family with the methods of distribution characteristic of a family. But the family is the embodiment of the private character of appropriation of the products of labour. In its modern form, it is itself a product of the development of private property. Insofar as the family carries out the production and reproduction of human beings as labour power (the human being as an economic agent of a society of commodity production), it is wholly embedded in, and entirely reproduced according to, the laws of motion of private property. Insofar as it reproduces all its members and serves as the basic consumer “unit”, it is wholly and completely determined by capitalist production. Economically, this is presupposed both in the form of wages and in the form of profit. Moreover, the modern family as an economic unit not only does not exclude but directly presupposes the exploitation of man by man, which is incompatible with communism. Therefore, family property cannot be regarded as the basis of social property—the property characteristic of communism.

And still less can communist property, which might serve as a nodal point in the movement towards communism, be found in all sorts of isolated communities unconnected with social production (for example, religious communes based on a natural economy).

But where, then, can we see the common property characteristic of communism? Clearly, only where there is free human activity in its directly collective form. Such activity, viewed from the standpoint of the use and disposal of resources, from the standpoint of appropriation by means of a social form, is precisely the actual existence of common property.

In the previous chapter we gave only a few examples of such communist activity in socialist and capitalist society. If we consider these and other examples from the standpoint of the exercise of will in relation to the subject of labour, the means of labour and the product, we are dealing with forms of social property and the social character of appropriation.

For example, a publicly accessible database, by its very nature, can no longer be anything other than social property, to say nothing of a common system of economic management. The same applies to free software and to the various services that presuppose joint production and use. There are also resource-sharing servers that would be inefficient and unprofitable to operate on the basis of private property—for example, shared libraries and repositories for films and music, which presuppose that participants both receive and distribute products on these services (although the example of software and databases is better, since it is more a matter of production than of consumption).

It is clear from the above that social property arises from the real movement of production even where there is as yet no formal socialisation of the means of production, under the conditions of capitalism. Formal socialisation—the seizure of political power by the proletariat and its taking possession of the means of production—will open up enormous scope for the development of real social property. Note that wherever there is real, and not merely formal, social property, the ground disappears for attempts to appropriate a portion of social resources for oneself, or, conversely, to squander them and treat them carelessly as though they belonged to “no one”.

Under formal socialisation, some external regulation is required—at least in the form of moral consciousness (in reality, in the form of law and a system of coercion that ensures its enforcement)—in relation to the means, subject and products of labour, so as to rule out such things. Under real socialisation, the need for such regulation falls away insofar as a culture of relating to the means of production becomes a simple condition of their use.
\end{document}
